\documentclass[UKenglish,usenames,dvipsnames,svgnames,table,aspectratio=169,mathserif]{beamer}

\mode<presentation> {

%\usetheme{default}
\usetheme{Madrid}

\setbeamertemplate{footline} % To remove the footer line in all slides uncomment this line

\setbeamertemplate{navigation symbols}{} % To remove the navigation symbols from the bottom of all slides uncomment this line
}

\usepackage{graphicx} % Allows including images
\usepackage{booktabs} % Allows the use of \toprule, \midrule and \bottomrule in tables
\usepackage{hyperref}
\usepackage{apacite}
\usepackage{babel}
\usepackage{fancyvrb}
\usepackage{color}
\usepackage{alltt}
\usepackage{listings}
\usepackage{framed}
\usepackage{courier}
\usepackage{minted}
\usepackage{epstopdf}
\usepackage{xifthen}
\usepackage[utf8]{inputenc}
\usepackage[T1]{fontenc}
\usepackage{textcomp}
\usepackage{gensymb}
\usepackage{svg}
\usepackage{pdfpages}
\usepackage{isodate}

\hypersetup{colorlinks=false}

\setbeamertemplate{bibliography entry title}{}
\setbeamertemplate{bibliography entry location}{}
\setbeamertemplate{bibliography entry note}{}
\setbeamertemplate{itemize items}[circle]
\setbeamertemplate{enumerate items}[circle]
\beamertemplatenavigationsymbolsempty
\setbeamertemplate{footline}{}


\newminted{haskell}{}
\newminted{scheme}{}

\lstdefinelanguage{haskell}{
  morekeywords={class,instance,where,do,data,newtype,default,deriving,module,import,qualified,as},
  otherkeywords={<-},
  sensitive=true,
  morecomment=[l]{--},
  morecomment=[n]{\{-}{-\}},
  morestring=[b]",
  morestring=[b]"""
}

\lstnewenvironment{code}{
  \lstset{
    language=haskell,
    basicstyle=\small\ttfamily,
    stringstyle=\color{olive}\ttfamily,
    escapeinside={*@}{@*},
    showstringspaces=false
  }
}{}

\definecolor{g}{RGB}{0,100,0}
\newcommand{\highlight}[1]{\colorbox{yellow}{#1}}
\newcommand{\nega}[1]{\colorbox{yellow}{#1}}
\newcommand{\posi}[1]{\colorbox{green}{#1}}
\newcommand{\nl}{\vspace{\baselineskip}}
\newcommand{\pnl}{\pause \nl}

\graphicspath{{images/}}

\newcommand{\textslide}[1]{{
\begin{frame}
\begin{center}

#1

\end{center}
\end{frame}
}}

\newcommand{\textslideleft}[1]{{
\begin{frame}

#1

\end{frame}
}}

\newcommand{\codeslide}[1]{{
\begin{frame}[fragile]
\begin{haskellcode}
#1
\end{haskellcode}
\end{frame}
}}


\newcommand{\imageslide}[2][1]{{
\begin{frame}\begin{center}
\includegraphics[scale=#1]{#2}
\end{center}\end{frame}
}}

\newcommand{\imageslideleft}[2][1]{{
\begin{frame}
\includegraphics[scale=#1]{#2}
\end{frame}
}}

\newcommand{\imagetextslide}[3][1]{{
\begin{frame}\begin{center}

{#3}

\includegraphics[scale=#1]{#2}
\end{center}\end{frame}
}}

\newcommand{\svgslide}[1]{{
\begin{frame}
\begin{center}
\includesvg{diagrams/#1}
\end{center}
\end{frame}
}}

\definecolor{bgc}{RGB}{255, 255, 255}
\setbeamercolor{background canvas}{bg=bgc}


%%----------------------------------------------------------------------------------------
%	TITLE PAGE
%----------------------------------------------------------------------------------------

\title[Education]{Functional Programming in Education}
\titlegraphic{\includegraphics[scale=0.2]{data61.eps}}
\author{George Wilson}
\institute[]
{
Data61/CSIRO\\
\medskip
\href{george.wilson@data61.csiro.au}{george.wilson@data61.csiro.au}
}

\selectlanguage{UKenglish}
\date{\printdate{2018-05-22}}

\begin{document}

%%%%%
%%%%% Intro section
%%%%%

\begin{frame}
\titlepage
\end{frame}


%%%%% SICP
\begin{frame}
\centering
\includegraphics{sicp.jpg}
\end{frame}


\begin{frame}[fragile]
\begin{schemecode}
(define (map proc items)
  (if (null? items)
      nil
      (cons (proc (car items))
            (map proc (cdr items)))))
\end{schemecode}
\end{frame}


\begin{frame}
\centering
\begin{tabular}{|l|l|}
\hline
       & Content \\
\hline\hline
Week 1 & Basic expressions  \\
\hline
Week 2 & procedure declarations \\
\hline
Week 3 & if-statement \\
\hline
Week 4 & while-statement \\
\hline
Week 5 & for-statement \\
\hline
\ldots & \\
\end{tabular}
\end{frame}


\begin{frame}[fragile]
\begin{schemecode}
(define (abs x)
  (cond ((> x 0) x)
        ((= x 0) x)
        ((< x 0) (- x))))
\end{schemecode}
\end{frame}


\begin{frame}
\centering
\includegraphics[scale=0.8]{sicp-factorial.png}
\end{frame}


\begin{frame}
\centering
\includegraphics{sicp-box-and-pointer.png}
\end{frame}


%%%%% Wadler
\begin{frame}
\centering
\includegraphics[scale=0.5]{wadler-title.png}
\end{frame}


\begin{frame}
\centering
\includegraphics[scale=0.55]{wadler-abstract.png}
% TODO highlight bits?
\end{frame}


\begin{frame}[fragile]
\begin{Verbatim}
(define (map proc items)
  (if (null? items)
      nil
      (cons (proc (car items))
            (map proc (cdr items)))))
\end{Verbatim}
\nl

\begin{haskellcode}
map :: (a -> b) -> [a] -> [b]
map f l =
  case l of
    []       -> []
    (x : xs) -> f x : map f xs
\end{haskellcode}
\end{frame}


\begin{frame}
\huge \centering
Thanks for listening!
\end{frame}


%\begin{frame}
%\frametitle{Attributions}
%\begin{itemize}
%\item
%  Structure and Interpretation of Computer Programs \\
%  Harold Abelson and Gerald Jay Sussman with Julie Sussman \\
%  Attribution-ShareAlike 4.0 International (CC BY-SA 4.0)
%\item
%  How to Design Programs \\
%  Matthias Felleisen, Robert Bruce Findler, Matthew Flatt, Shriram Krishnamurthi \\
%  Attribution-NonCommercial-NoDerivs 2.0 Generic (CC BY-NC-ND 2.0)
%\end{itemize}
%\end{frame}


\end{document}

